\documentclass[11pt]{article}

\usepackage[margin=1in]{geometry}
\usepackage[T1]{fontenc}
\usepackage[utf8]{inputenc}
\usepackage{lmodern}
\usepackage{microtype}
\usepackage{booktabs}
\usepackage{array}
\usepackage{graphicx}
\usepackage{url}
\usepackage[hidelinks]{hyperref}
\usepackage{orcidlink}

\newcommand{\code}[1]{\texttt{#1}}
% [FILL] markers flag places where the author must insert a real, verified value
% or write original prose. Do not submit with any [FILL] remaining.
\newcommand{\FILL}[1]{\textbf{[FILL: #1]}}

\title{Evidence-Driven Validation of a Hybrid Network Intrusion\\
Detection System: From Offline Replay to Prepared Live Capture}

\author{Azhad Shahzad Shaik\,\orcidlink{0009-0009-6450-5837}\\
\small \href{mailto:shaikazhadshahzad@gmail.com}{shaikazhadshahzad@gmail.com} \quad
\small \url{https://github.com/s-shahzad/Universal-NIDS}}

\date{\FILL{month year}}

\begin{document}
\maketitle

\begin{abstract}
Most network intrusion detection (NIDS) research reports offline accuracy on a
labelled dataset and stops there. We argue that for a hybrid detector intended
for real deployment, the validation process is itself a contribution. We present
a hybrid NIDS that combines deterministic signature rules, statistical anomaly
detection, a supervised ensemble (random forest, extra-trees, histogram gradient
boosting, and XGBoost), an optional unsupervised path (isolation forest plus a
shallow autoencoder), and a fusion layer that adjudicates the component outputs
into a single alert decision. Alongside the detector we describe a layered,
reproducible validation method: a deterministic test suite with an enforced
coverage floor, repeatable offline PCAP-replay scenarios, and prepared-environment
runs on live capture that measure benign false positives, queue-loss accounting,
restart recovery, and sustained-soak behaviour. On a held-out test split of the
labelled traffic corpus the supervised ensemble reaches
\textbf{99.5\% accuracy and 0.996 weighted $F_1$} across five classes, while the
per-class results expose a genuine weakness on the rarest class that offline
aggregates hide. The prepared-environment evidence shows the live path, loss
accounting, operator workflows, and a completed six-hour soak (87{,}533 flows, 0
alerts, $\approx$400\,MiB peak memory) all function as intended. We still
characterise the system as a controlled pre-deployment candidate rather than a
production system, because benign adjudication and platform coverage remain
narrow, and we release the code, datasets, and evidence records so the validation
is reproducible.
\end{abstract}

\section{Introduction}

Network intrusion detection is hard for two reasons that are usually addressed
separately. The first is \emph{detection quality}: distinguishing malicious from
benign traffic. The second is \emph{operational credibility}: showing that a
detector behaves acceptably on live traffic, under benign drift, and through the
maintenance actions an operator actually performs. Benchmark-oriented NIDS papers
answer the first and rarely the second; a high offline $F_1$ says little about
false-positive burden on live capture or about whether the system survives a
restart without losing continuity.

This paper treats a hybrid NIDS \emph{and its validation process} as one research
object. The detector combines five detection paths and a fusion layer
(\S\ref{sec:arch}). The validation method (\S\ref{sec:results}) is layered:
deterministic tests, repeatable offline replay, and prepared-environment live
runs that measure the things offline metrics cannot, benign false positives,
packet-queue loss, restart recovery, and soak stability. Our position is that
this second half is where deployment confidence is earned or lost.

\paragraph{Contributions.}
\begin{enumerate}
  \item A hybrid NIDS architecture that retains component-level evidence for every
        decision, so an analyst can see why an alert did or did not fire
        (\S\ref{sec:arch}, \S\ref{sec:method}).
  \item A layered, reproducible validation method spanning code tests, offline
        replay, and prepared-environment live capture (\S\ref{sec:results}).
  \item An honest evaluation that reports both the strong aggregate offline
        result and the per-class and live-capture weaknesses it hides
        (\S\ref{sec:results}, \S\ref{sec:limits}).
  \item A released harness, labelled dataset split, and evidence records for
        reproducibility.
\end{enumerate}

\section{Related Work}
\label{sec:related}

We group prior NIDS work into four families and position ours against them.

\textbf{Signature and rule engines} such as Suricata and Zeek are precise and
interpretable but weak on behaviour not covered by a rule; our signature engine
plays the same role and is deliberately only one of five paths.

\textbf{Supervised-ML IDS on labelled datasets} dominate the literature and report
strong offline scores. Much of this work is evaluated on the KDD~Cup/NSL-KDD
family~\cite{nslkdd}, whose \code{dos}/\code{probe}/\code{r2l}/\code{u2r} attack
taxonomy we adopt for labels, and on the flow-based CIC-IDS2017
corpus~\cite{cicids2017}. Ensemble and feature-selection classifiers push
aggregate accuracy high on these sets~\cite{ensembleids}, but the community has
also documented that headline numbers on these datasets can be fragile: label and
processing errors in CIC-IDS2017 materially change reported
performance~\cite{cicidstrouble}. This is exactly why we report per-class results
and treat offline accuracy as one input, not the verdict.

\textbf{Anomaly and statistical IDS} improve sensitivity to novel behaviour but
are prone to benign-drift false positives; our unsupervised path is gated
precisely to contain that failure mode.

\textbf{Hybrid and validation-oriented work} is closest to ours. Our own prior
work applied supervised ML to denial-of-service detection on IoT EV-charging
infrastructure using a CIC-IDS-derived binary corpus~\cite{shaik2025ocpp}; the
present work is distinct in three ways: it is a five-path hybrid rather than a
single classifier, it evaluates a five-class taxonomy rather than binary
DoS/benign, and its central contribution is the prepared-environment validation
method rather than the offline score. We build against operational guidance from
NIST SP~800-115, NIST SP~800-61r3, MITRE ATT\&CK, and CISA logging/SIEM
guidance~\cite{nist800115,nist80061,attack,cisalog}.

\section{System Architecture}
\label{sec:arch}

The system ingests from live NIC capture, offline PCAP replay, optional
Suricata/Zeek JSON adapters, and static artifact intake, then normalises all
inputs into a common flow-event schema before feature extraction. Normalised
features fan out to the detection engines and a fusion layer, and results persist
to SQLite and JSONL with an analyst dashboard and reporting on top.

\begin{figure}[h]
\centering
\FILL{include architecture diagram: export thesis/diagrams/system\_architecture.mmd
to PDF (mermaid-cli) into paper/figures/architecture.pdf, then
\textbackslash includegraphics here}
\caption{End-to-end architecture: ingest, normalise, five detection paths,
fusion, suppression, persistence, and reporting.}
\label{fig:arch}
\end{figure}

Artifact analysis is static only; submitted files are never executed.

\section{Detection Methodology}
\label{sec:method}

\begin{itemize}
  \item \textbf{Signature engine.} YAML-defined rules for known behaviour; produces
        reviewable, deterministic alerts.
  \item \textbf{Statistical anomaly engine.} EWMA and z-score style thresholds with
        burst detection over flow aggregates.
  \item \textbf{Supervised ensemble.} Random forest, extra-trees, histogram
        gradient boosting, and XGBoost (when available), weighted into a single
        probability vector. \FILL{state ensemble weighting / voting scheme from
        src/NIDS/ml/supervised\_ensemble.py in one precise sentence.}
  \item \textbf{Unsupervised path (optional).} Isolation forest plus a shallow
        autoencoder with warmup calibration; gated to require at least two active
        components before it can raise an alert, a false-positive control
        introduced during tuning (\S\ref{sec:results}).
  \item \textbf{Fusion.} Combines component scores into a benign / informational /
        alert decision while retaining each component's contribution as evidence.
\end{itemize}

\section{Datasets and Experimental Setup}
\label{sec:data}

\paragraph{Labelled corpus.} The supervised ensemble is trained and evaluated on a
labelled traffic corpus of \textbf{25{,}000 flows} across five classes
(\code{dos}, \code{normal}, \code{probe}, \code{r2l}, \code{u2r}), split
\textbf{18{,}750 train / 6{,}250 held-out test}. The label taxonomy is the
KDD-family \code{dos}/\code{normal}/\code{probe}/\code{r2l}/\code{u2r}
scheme~\cite{nslkdd}, but the features are \emph{not} the classic 41 KDD
attributes: the model consumes this system's own 15-column flow schema (packet
length, ports, protocol and TCP-flag indicators, per-window destination-port and
-host counts, and DNS/HTTP/TLS presence flags) from
\code{src/NIDS/ml/featureset.py}. The labels follow KDD conventions while the
feature space is the deployed detector's own. \FILL{state precisely how the
labelled training flows were produced or sourced (own labelled capture vs.\
re-mapped public corpus) and cite it; this must be exact before submission.}
We report the held-out split as the
primary result; a second evaluation over the full 25{,}000 flows is reported for
comparison and is \emph{not} a held-out number. The class distribution is heavily
imbalanced: \code{u2r} is so rare that it has no representative in the held-out
test split, and \code{probe} and \code{r2l} have only tens of samples. We keep
this visible rather than smoothing it into a single accuracy figure.

\paragraph{Live and offline scenarios.} Offline PCAP-replay scenarios
(\code{LAB-SCN-001..005}: port scan, HTTP brute force, flood/burst, mixed
benign+malicious, artifact+network correlation) and \textbf{thirteen}
prepared-environment live scenarios (\code{PREP-ENV-001..013}: live capture,
malformed-packet handling, queue-loss accounting, restart recovery, benign soak,
full-duration soak, operator workflows, and live suppression) are executed
through released scripts and indexed as manifests. At draft time the index
records 36 runs with all thirteen prepared-environment scenarios passing.

\section{Results}
\label{sec:results}

\subsection{Software validation}

\begin{table}[h]
\centering
\begin{tabular}{lr}
\toprule
Measure & Value \\
\midrule
Tests collected (public CI scope) & 245 \\
Passing (default suite) & 229 \\
Deselected (lab/environment/live markers) & 16 \\
Active pytest warnings & 0 \\
Coverage (line) & 75.83\% \\
Enforced coverage floor & 72\% \\
\bottomrule
\end{tabular}
\caption{Deterministic test-suite result at draft time (public CI scope). CI is
green on all three workflows.}
\label{tab:tests}
\end{table}

\subsection{Supervised ensemble, held-out test split}

\begin{table}[h]
\centering
\begin{tabular}{lrrrr}
\toprule
Class & Precision & Recall & $F_1$ & Support \\
\midrule
dos     & 0.999 & 0.996 & 0.997 & 4941 \\
normal  & 0.998 & 0.995 & 0.996 & 1243 \\
probe   & 0.725 & 0.980 & 0.833 & 51 \\
r2l     & 0.933 & 0.933 & 0.933 & 15 \\
u2r     & \multicolumn{4}{c}{\emph{no samples in this test split}} \\
\midrule
\textbf{Weighted avg} & \textbf{0.996} & \textbf{0.995} & \textbf{0.996} & 6250 \\
Accuracy & \multicolumn{4}{r}{\textbf{0.9954}} \\
\bottomrule
\end{tabular}
\caption{Supervised ensemble on the 6{,}250-flow held-out test split. Aggregate
accuracy is 99.5\%, but the \code{probe} and \code{r2l} classes, with 51 and 15
test samples, are materially weaker (precision 0.72 on \code{probe}), and the
\code{u2r} class has no samples in this split at all, so nothing can be claimed
about it. These are the honest costs that a single weighted aggregate hides.}
\label{tab:mlheld}
\end{table}

The aggregate is strong, but Table~\ref{tab:mlheld} is deliberately reported
per-class: the rare classes are where the ensemble is weakest, and any deployment
claim must account for the small support behind their scores. We treat these as
offline-evaluation numbers, not live operational accuracy.

\subsection{Prepared-environment (live) results}

\begin{table}[h]
\centering
\begin{tabular}{lp{8.8cm}}
\toprule
Scenario & Result \\
\midrule
Queue-loss accounting (PREP-ENV-003) & Under a deliberate DNS flood, 8{,}127 packets received, 23 processed, 8{,}100 dropped (99.67\% loss), queue-depth peak 1, 0 false alerts. The point is that loss is \emph{accounted}, not that it is avoided. \\
Benign soak, tuned (PREP-ENV-005) & 1{,}404 flows over 180\,s, 0 alerts; the two unsupervised false positives seen before tuning dropped to 0. \\
Restart recovery (PREP-ENV-006) & Two-phase restart, 144 flows, continuity preserved; reload latency $\approx$13.2\,s. \\
Full-duration soak (PREP-ENV-007) & \textbf{Completed} 6-hour run (21{,}600\,s), 87{,}533 flows, 0 alerts, 0 dropped packets, peak RSS 409{,}444\,KiB ($\approx$400\,MiB), peak CPU 107\%, restart latency 13.335\,s. \\
Operator workflows & Rule refresh (PREP-ENV-008): 321 flows, 1 alert. Model swap (PREP-ENV-009): 330 flows, 2 alerts. Config override (PREP-ENV-010): 394 flows, 0 alerts. Each reload $\approx$13.2--13.3\,s. \\
\bottomrule
\end{tabular}
\caption{Prepared-environment live-capture evidence, from
\code{NIDS\_TestLab/reports/prepared\_env\_validation\_index.json}
(2026-03-14, 36 recorded runs, all 13 scenarios passing). Values frozen at draft
time.}
\label{tab:prepenv}
\end{table}

\subsection{False-positive tuning}

Early prepared-environment runs raised unsupervised-only alerts on benign status
polling and routine DNS. We added a control requiring at least two active
unsupervised components before an alert can fire; the tuned rerun cleared the
exercised benign sample (\FILL{PREP-ENV-005: flows, 0 alerts}). We claim this only
for the exercised sample; broader benign corpora are future work.

\section{Limitations}
\label{sec:limits}

\begin{itemize}
  \item \textbf{Single soak run.} The six-hour soak completed cleanly, but it is
        one run at one traffic profile; we do not generalise it to arbitrary
        deployment loads or durations.
  \item \textbf{Rare-class detection.} \code{probe}/\code{r2l} have small test
        support and weaker scores, and \code{u2r} is absent from the held-out
        split (Table~\ref{tab:mlheld}).
  \item \textbf{Benign adjudication is narrow.} Tuning is validated on the
        exercised sample, not a broad benign corpus.
  \item \textbf{Maintenance.} Operator workflows are restart-based, implying an
        $\approx$13-second blind window per reload; not zero-downtime.
  \item \textbf{Platform scope.} \FILL{state validated OS/interface scope honestly
        (Windows first-class lab host, Linux runtime target per the repo docs);
        confirm before submission.}
\end{itemize}

\section{Conclusion}

We presented a hybrid NIDS whose contribution is as much its layered, reproducible
validation as its detection stack. The supervised ensemble is strong in aggregate
(99.5\% held-out accuracy) but honestly weaker on rare classes, and the
prepared-environment evidence, including a completed six-hour live soak, supports
a \emph{controlled pre-deployment candidate} verdict rather than production
readiness. We release the harness,
dataset split, and evidence records so the result is reproducible.

\begin{thebibliography}{9}
\bibitem{nist800115} NIST. \emph{SP 800-115: Technical Guide to Information
Security Testing and Assessment.} \url{https://csrc.nist.gov/pubs/sp/800/115/final}
\bibitem{nist80061} NIST. \emph{SP 800-61r3: Incident Response Recommendations.}
\url{https://csrc.nist.gov/pubs/sp/800/61/r3/final}
\bibitem{attack} MITRE. \emph{ATT\&CK Enterprise Matrix.}
\url{https://attack.mitre.org/matrices/enterprise/}
\bibitem{cisalog} CISA et al. \emph{Best Practices for Event Logging and Threat
Detection.} \url{https://www.cisa.gov/resources-tools/resources/best-practices-event-logging-and-threat-detection}

\bibitem{nslkdd} M. Tavallaee, E. Bagheri, W. Lu, and A. A. Ghorbani.
\newblock A detailed analysis of the KDD CUP 99 data set.
\newblock In \emph{IEEE Symposium on Computational Intelligence for Security and
Defense Applications (CISDA)}, 2009. (NSL-KDD dataset.)

\bibitem{cicids2017} I. Sharafaldin, A. H. Lashkari, and A. A. Ghorbani.
\newblock Toward generating a new intrusion detection dataset and intrusion
traffic characterization.
\newblock In \emph{4th Intl. Conf. on Information Systems Security and Privacy
(ICISSP)}, 2018. (CIC-IDS2017 dataset.)

\bibitem{ensembleids} Y. Zhou, G. Cheng, S. Jiang, and M. Dai.
\newblock Building an efficient intrusion detection system based on feature
selection and ensemble classifier.
\newblock \emph{Computer Networks}, 174:107247, 2020. Preprint
\url{https://arxiv.org/abs/1904.01352}.

\bibitem{cicidstrouble} G. Engelen, V. Rimmer, and W. Joosen.
\newblock Troubleshooting an intrusion detection dataset: the CICIDS2017 case
study.
\newblock In \emph{IEEE Security and Privacy Workshops (WTMC)}, 2021.

\bibitem{shaik2025ocpp} A. S. Shaik, M. Gutla, A. P. R. Kagitala, and S. Bhatia.
\newblock Plugged-in and protected: Leveraging machine learning to secure
IoT-based electric vehicle charging stations from denial-of-service threats.
\newblock In \emph{IEEE Global Conference on Artificial Intelligence and Internet
of Things (GCAIoT)}, 2025. DOI: 10.1109/GCAIoT68269.2025.11275540.

\end{thebibliography}

\end{document}
